\documentclass[12pt]{article}
\usepackage{amsmath,amssymb,amsthm}
\usepackage{geometry}
\usepackage{booktabs}
\usepackage{hyperref}
\usepackage{microtype}
\geometry{margin=1in}

\newtheorem{proposition}{Proposition}
\newtheorem{theorem}{Theorem}
\newtheorem{corollary}{Corollary}
\newtheorem{definition}{Definition}

\newcommand{\taubar}{\bar{\tau}}
\newcommand{\Phifield}{\Phi}

\title{\textbf{Module 3: Shell Hierarchy and the Nested Fold Field}\\[4pt]
\large From Microscopic Shell Formation to Cosmological Structure}

\author{Sean Sowden\\
\small Prime Fold Theory Project\\[4pt]
\small Mathematical assistance: Claude Sonnet~4.6 --- ``Green'' (Anthropic, 2026)}

\date{April 2026}

\begin{document}
\maketitle

\begin{abstract}
Module~1 established that three-dimensional space emerges when regions of the
Fold Field close upon themselves, forming shells. This paper extends that result
upward in scale: the same shell-formation dynamics that produced 3D geometry
operate at all scales, organizing matter into a nested hierarchy of gravitationally
bound shells. We derive the fold-native shell boundary condition --- the
Lagrange surface where one fold field yields to another --- and show how the
background structural mass term $B$ transforms across shell boundaries.
The universe itself is the outermost shell, closed during the 2D$\to$3D
transition. The differential $B_{\mathrm{boundary}} - B_{\mathrm{interior}}$,
built up as structure forms within the cosmological shell, is identified as a
candidate mechanism for dark energy. The timing is natural: the dark sector
cannot appear before shell closure, which is why it arrives late.
All derivations use fold primitives $(\Phi, \taubar, \kappa)$ and the
results of Papers~1--4 in the PFT series.
\end{abstract}

%──────────────────────────────────────────────────────────────────────
\section{Introduction}

The dimensional ladder of Module~1 runs upward from the Prime Substrate:
\begin{equation}
  \text{0D} \;\to\; \text{1D} \;\to\; \text{2D (Fold Field)} \;\to\;
  \text{3D (Shells)}.
\end{equation}
Shells are the mechanism by which dimensionality crystallizes: a region of
the Fold Field closes upon itself, acquiring interior/exterior structure
and quasi-radial layering. Once shells exist, gravity follows
(Paper~3~\cite{gravity}) and gauge structure organizes
(Papers~1--2~\cite{gauge,attractor}).

That account stops at the Planck scale. It explains how 3D space emerged,
but not how matter organized itself into stars, planets, and galaxies within
that space. This paper argues the gap is not a new mechanism but the same
one: shell formation continues operating at every scale above the Planck
length, producing a nested hierarchy that extends from subatomic to
cosmological scales.

\subsection{The Central Claim}

\begin{quote}
\textit{We hypothesize that the fold field is scale-invariant in
its shell-formation dynamics: a region of the fold field closes
into a shell whenever its local $g_{\mathrm{local}}(r) >
g_{\mathrm{external}}(r)$, regardless of scale, and the shell
boundary is the surface where these fields cross.
The boundary condition is derived; scale invariance of the
full dynamics is a hypothesis supported by the absence of
explicit scale dependence, not yet a proven theorem.}
\end{quote}

The boundary condition itself follows from the inverse-square law
(Paper~3) and the superposition principle (which holds exactly in
the static limit) --- these are results, not postulates.
The claim that this same condition applies at all scales is the
hypothesis. The novelty is in tracing its consequences: a nested
$B$-hierarchy, screening of the cosmological background, and a
natural candidate for dark energy.

\subsection{What This Paper Does Not Claim}

\begin{itemize}
  \item We do not derive the magnitude of dark energy from fold primitives.
        That requires a fold cosmology not yet developed.
  \item We do not prove scale invariance of shell formation.
        We derive the boundary condition and note that it has no
        explicit scale dependence; we flag this as a hypothesis
        requiring further verification.
  \item We do not identify the dark matter candidate precisely.
        We note the natural slot for a high-$\kappa$ fold sector and
        defer its development.
\end{itemize}

%──────────────────────────────────────────────────────────────────────
\section{Shell Formation: From Planck Scale to Macroscopic}

\subsection{The Module~1 Shell}

Module~1 defines a shell as a closed region of the Fold Field with:
\begin{enumerate}
  \item A defined interior/exterior boundary
  \item Quasi-radial layering (the fold density $\Phifield$ stratifies
        by distance from the shell center)
  \item Net skew $\mathcal{S} = \nabla\cdot g \neq 0$ at the center
        (a mass source in the sense of Paper~3)
\end{enumerate}
Three-dimensional space itself is such a shell --- the outermost one,
formed during the 2D$\to$3D transition driven by cost-of-connectivity
dynamics (spectral dimension $d_s \approx 3.1$, confirmed numerically).

\subsection{The Macroscopic Shell}

At scales above the Planck length, the fold field has already crystallized
into 3D geometry. A macroscopic object (planet, star, galaxy) is a region
where the local fold density $\Phifield$ and structural mass $m = \Phifield\taubar$
are elevated above the background. The gravitational field of this object:
\begin{equation}
  g_{\mathrm{local}}(r) = \frac{A_{\mathrm{local}}}{r^2}
\end{equation}
competes with the gravitational field of any containing object:
\begin{equation}
  g_{\mathrm{ext}}(r) = \frac{A_{\mathrm{ext}}}{(d-r)^2}
\end{equation}
where $d$ is the distance to the external source.

\begin{definition}[Shell Boundary]
  The shell boundary of a fold object with structural mass parameter
  $A_{\mathrm{local}}$ within an external field with parameter
  $A_{\mathrm{ext}}$ at distance $d$ is the surface $r = r^*$ satisfying:
  \begin{equation}
    \frac{A_{\mathrm{local}}}{r^{*2}} = \frac{A_{\mathrm{ext}}}{(d-r^*)^2}.
    \label{eq:boundary}
  \end{equation}
  Inside $r^*$: local fold dynamics dominate.
  Outside $r^*$: external fold dynamics dominate.
\end{definition}

This is the fold-native form of the Hill sphere / Lagrange surface.
It is not a postulate. Module~2.3 \cite{mod23} (Section~4) derives it analytically
from $1/r^2$ and superposition, and confirms it numerically on the
fold lattice with correlation $r=0.994$ across five mass ratios
(errors $<6\%$, consistent with lattice discretization).
The closed-form solution is:
\begin{equation}
  r^* = \frac{d}{1 + \sqrt{A_{\mathrm{ext}}/A_{\mathrm{local}}}},
  \label{eq:hill}
\end{equation}
which follows from setting $g_{\mathrm{local}}(r^*) = g_{\mathrm{ext}}(r^*)$
and solving. The larger source dominates the larger volume ---
``bias to the larger'' is a mathematical consequence of the fold
field, not an additional assumption.

\begin{proposition}[Scale Independence of Boundary Condition]
  Equation~\eqref{eq:boundary} contains no explicit length scale.
  The shell boundary exists at any scale where a local fold field
  competes with an external one, from nuclear to cosmological.
\end{proposition}

\textbf{Nondimensionalization argument.}
To verify that scale invariance is not merely absence of explicit
scale but a structural property, we nondimensionalize the fold
field equations. Let $r_0$ be any reference length, $m_0$ any
reference mass, $t_0 = r_0/c$ the associated time. Define
dimensionless variables:
\begin{equation}
  \tilde{r} = r/r_0, \quad
  \tilde{m} = m/m_0, \quad
  \tilde{t} = t/t_0, \quad
  \tilde{\taubar} = \taubar/t_0, \quad
  \tilde{\kappa} = \kappa/t_0.
\end{equation}
The fold transport equation, accumulation equation, and shell
boundary condition~\eqref{eq:boundary} all transform to equations
in tilded variables with \emph{identical form} and no residual
dependence on $r_0$, $m_0$, or $t_0$.
The ratio $\tilde{A}/\tilde{A}_{\mathrm{ext}}$ determining $r^*/d$
is invariant under this rescaling.

This confirms that the shell boundary condition is not merely
lacking an explicit scale --- it is genuinely scale-free.
The same condition applies at Planck scale, planetary scale,
and cosmological scale because no combination of $r_0$, $m_0$, $t_0$
appears in the rescaled equations.

\textit{Caveat:} scale invariance of the boundary condition does
not guarantee scale invariance of the full dynamical solution.
Boundary conditions and initial conditions may break the symmetry
at specific scales. This is Open Problem~2 (Section~\ref{sec:open}).

\subsection{The Rocket Thought Experiment}

A rocket launching from Earth illustrates the shell hierarchy concretely.
Inside Earth's Hill sphere ($r < r^*_{\mathrm{Earth}}$): the rocket's
fold field contributes to Earth's $M_{\mathrm{structure}}$ and is screened
by Earth's shell. The rocket has no independent shell.

As the rocket crosses $r^*_{\mathrm{Earth}}$: the local/external field
balance reverses. The rocket's fold field now dominates locally. A new
shell boundary forms around the rocket, with the rocket's $A_{\mathrm{rocket}}$
competing against Earth's field at the new distance. The rocket has claimed
its own shell.

The boundary is not determined by the rocket's decision or velocity.
It is determined by the ratio $A_{\mathrm{local}}/A_{\mathrm{ext}}$
--- bias toward the larger. A more massive object always wins within its
Hill sphere.

%──────────────────────────────────────────────────────────────────────
\section{The $B$-Hierarchy: Screening and the Nested Background}

\subsection{$B$ in the General Solution}

The static fold field equation $\nabla^2 m = 0$ has general solution:
\begin{equation}
  m(r) = \frac{A}{r} + B.
\end{equation}
$B$ is spatially uniform within its domain, drops out of the local force
law ($\nabla B = 0$), and is invisible to local dynamics. In Paper~3
this was noted as a cosmological constant analog. The shell hierarchy
gives it a precise meaning.

\subsection{$B$ Transforms Across Shell Boundaries}

Inside a shell at scale $\ell$, the relevant background is not the
cosmological $B_{\mathrm{universe}}$ but the local $B_\ell$ set by
the shell's own T3 conservation:
\begin{equation}
  B_\ell = \frac{M_{\mathrm{total},\ell} - M_{\mathrm{structure},\ell}}{V_\ell}.
  \label{eq:B_local}
\end{equation}

\textbf{Derivation of screening from Definition~2.1.}
By Definition~2.1, inside the shell boundary $r < r^*$, the local
fold field $g_{\mathrm{local}}(r)$ dominates over the external field
$g_{\mathrm{ext}}(r)$.
The general solution $m(r) = A/r + B$ applies separately in each
domain: inside the shell, $m_{\mathrm{in}}(r) = A_{\mathrm{local}}/r + B_\ell$;
outside, $m_{\mathrm{out}}(r) = A_{\mathrm{ext}}/r + B_{\mathrm{universe}}$.
Continuity of $m$ at $r^*$ requires:
\begin{equation}
  \frac{A_{\mathrm{local}}}{r^*} + B_\ell =
  \frac{A_{\mathrm{ext}}}{r^*} + B_{\mathrm{universe}}.
\end{equation}
Since $g_{\mathrm{local}}(r^*) = g_{\mathrm{ext}}(r^*)$ at the boundary
(by Definition~2.1), the $A/r$ terms match and we require $B_\ell$
to absorb any residual --- it is set by T3 within the shell volume,
not inherited from the exterior.
The cosmological $B_{\mathrm{universe}}$ does not penetrate: it is
present at the boundary but is replaced by $B_\ell$ in the interior
domain, because T3 conservation within $V_\ell$ independently fixes
the local mass budget.
Screening is therefore a consequence of the boundary condition
(Definition~2.1) combined with T3 --- not of T3 alone.

\begin{proposition}[$B$-Screening]
  Within a fold shell of volume $V_\ell$, the effective background
  structural mass is $B_\ell$ given by Eq.~\eqref{eq:B_local},
  derived from the shell boundary condition (Definition~2.1) and
  T3 conservation.
  The cosmological background $B_{\mathrm{universe}}$ does not
  contribute to local dynamics inside the shell.
\end{proposition}

This explains why the cosmological constant does not affect planetary
orbits, solar system dynamics, or galactic rotation curves at small
radii --- the local shell screens it completely. Dark energy only acts
on objects that have exited all containing shells, i.e.\ objects in
pure Hubble flow.

\subsection{The Nested Hierarchy}

The full $B$-hierarchy is:
\begin{equation}
  B_{\mathrm{universe}} \supset B_{\mathrm{cluster}} \supset
  B_{\mathrm{galaxy}} \supset B_{\mathrm{solar}} \supset
  B_{\mathrm{planet}} \supset \cdots
\end{equation}
Each level screens the one above. At each boundary crossing,
Eq.~\eqref{eq:B_local} sets the new local $B$ from the enclosed
mass budget.

\textbf{Dynamical hierarchy formation.}
A new shell boundary forms at scale $L$ when the local fold
accumulation $\taubar$ reaches threshold $\kappa$ across a
region of that scale, triggering a collapse event that
concentrates structural mass and establishes an interior/exterior
contrast. The interior $B$ then decouples from the external field
via T3 conservation (Eq.~\ref{eq:B_local}), and a new screening
layer is established.

In plain language: \emph{a new shell forms when local tension
reaches its breaking point at a given scale, concentrating mass,
drawing a new boundary, and draining the background field inside it.}
The scale at which this occurs is set by the ratio $\kappa/\alpha$
and the local gradient $|\nabla m|$ --- the same parameters that
set the shell plateau $C = 1.014$ in Module~2.3.
Each level of the hierarchy is the fold dynamics solving the same
fixed-point problem at a different scale.

%──────────────────────────────────────────────────────────────────────
\section{The Cosmological Shell and Dark Energy}

\subsection{The Universe as the Outermost Shell}

The universe is the outermost fold shell --- the one that closed during
the 2D$\to$3D transition. Its formation is not an event in time within
the universe; it is the event that produced time (via the Planck tick,
Paper~4~\cite{planck_tick}) and 3D space (via cost-of-connectivity
dynamics). There is no containing shell above it.

Consequently, the universe's $B_{\mathrm{boundary}}$ has no external
reference to screen it. It is set by the initial conditions at shell
closure --- the total structural mass of the universe at the moment of
3D crystallization.

\subsection{The $B$-Differential and Dark Energy}

As structure forms within the cosmological shell, $M_{\mathrm{structure}}$
grows. By Eq.~\eqref{eq:B_local}:
\begin{equation}
  B_{\mathrm{interior}}(t) = \frac{M_{\mathrm{total}} -
  M_{\mathrm{structure}}(t)}{V(t)}.
\end{equation}
$B_{\mathrm{interior}}$ decreases as structure grows.
The shell boundary $B_{\mathrm{boundary}}$, set at the 2D$\to$3D
transition, does not decrease at the same rate --- it is anchored
to the closed-shell conditions.

The differential:
\begin{equation}
  \Delta B(t) = B_{\mathrm{boundary}} - B_{\mathrm{interior}}(t) \geq 0
  \label{eq:deltaB}
\end{equation}
grows monotonically as structure forms. In the HVAC analogy:
the interior is at lower structural-mass pressure than the boundary.
The boundary does not need to be negative in absolute terms ---
only higher than the interior. This is suction, not absolute negative
pressure.

\begin{proposition}[$\Delta B$ as Effective Negative Pressure]
  In a closed fold system (T3 conservation, shell boundary fixed),
  the differential $\Delta B = B_{\mathrm{boundary}} -
  B_{\mathrm{interior}} > 0$ acts as an effective outward pressure
  on the interior. It does not require $B < 0$ and does not violate T1.
  It requires only that $B_{\mathrm{boundary}}$ decreases more slowly
  than $B_{\mathrm{interior}}$ as structure forms.
\end{proposition}

\subsection{Why the Dark Sector Arrives Late}

The cosmological $\Delta B$ cannot exist before shell closure.
Before the 2D$\to$3D transition, there is no shell boundary, no
interior/exterior distinction, and no $B$-differential. The dark sector
--- whatever manifests $\Delta B$ as observable dark energy --- is
structurally forbidden from appearing before 3D space exists.

After shell closure, $\Delta B = 0$ initially (no structure has formed).
$\Delta B$ grows as $M_{\mathrm{structure}}$ grows. Dark energy becomes
dynamically significant when $\Delta B$ becomes comparable to the
local gravitational energy density --- which happens around the epoch
of matter-energy equality, consistent with observations ($z \approx 0.4$).

This resolves the coincidence problem naturally: dark energy and matter
are comparable today because dark energy is \emph{sourced by} matter
clustering. They are not independent quantities that happen to coincide.

\subsection{Equation of State}

The equation of state of $\Delta B$ depends on how $B_{\mathrm{boundary}}$
and $B_{\mathrm{interior}}$ evolve with the expansion scale factor $a$.

If $B_{\mathrm{boundary}} = \mathrm{const}$ (boundary frozen at
shell-closure conditions) and $B_{\mathrm{interior}} \propto a^{-3}$
(dilutes with volume):
\begin{equation}
  \Delta B(a) \approx B_{\mathrm{boundary}} - B_0 a^{-3}
  \xrightarrow{a \to \infty} B_{\mathrm{boundary}} = \mathrm{const}.
\end{equation}
At late times $\Delta B \to \mathrm{const}$, giving $w = -1$.
At early times $\Delta B \approx B_{\mathrm{boundary}} - B_0 a^{-3}$
is smaller and growing, giving $w > -1$ (quintessence-like).

The transition from $w > -1$ (early) to $w \to -1$ (late) is a
prediction: dark energy should appear to strengthen slightly over cosmic
time as $B_{\mathrm{interior}}$ depletes. Current observations
(DESI 2024 BAO results) are consistent with $w$ slightly above $-1$
and possibly time-varying, which is consistent with this picture.

\subsection{What This Is Not}

This is a structural identification, not a derivation of the dark energy
density. The magnitude of $B_{\mathrm{boundary}}$ --- why it corresponds
to $\sim 10^{-29}$ g/cm$^3$ --- requires a fold cosmology that connects
the Planck-scale shell closure conditions to the observed universe. That
cosmology does not yet exist.

%──────────────────────────────────────────────────────────────────────
\section{The Dark Matter Slot}

\subsection{Two Fold Sectors}

The fold framework currently treats all nodes identically: same $\Phi$,
same $\taubar$, same $\kappa$. A natural extension is two populations:

\begin{itemize}
  \item \textbf{Visible sector}: nodes with threshold $\kappa_v$,
        coupling to the electromagnetic fold cycle (Papers~1--2).
  \item \textbf{Dark sector}: nodes with threshold $\kappa_d \gg \kappa_v$,
        no electromagnetic coupling.
\end{itemize}

High-$\kappa$ nodes accumulate $\taubar$ longer before resetting.
They form deeper potential wells. They cluster more readily.
They do not participate in the color-neutrality attractor (Paper~2)
because they never reach the 3-state cycle threshold.

\subsection{Why It Arrives with Shell Closure}

Dark sector nodes are present in the fold field from the beginning ---
they are simply fold nodes with high $\kappa$. But they become
gravitationally significant only after shell closure, when the
inverse-square law (Paper~3) begins operating. Before 3D space exists,
there is no distance, no $1/r^2$, no gravitational clustering.

After shell closure, dark sector nodes cluster first (deeper wells,
longer accumulation, stronger local $g$) and then visible matter
falls into the resulting structure. This is the standard dark matter
picture --- but derived from fold dynamics rather than postulated.

\subsection{Honest Status}

This is a candidate identification. The two-sector fold field has not
been simulated. The coupling between sectors (through the shared
gravitational field) has not been derived. This is noted here as
the natural dark matter slot in the framework, to be developed in
future work.

%──────────────────────────────────────────────────────────────────────
\section{Quaoar as a Shell Boundary Test}

The Quaoar ring anomaly (Paper~3, Open Problem discussion) is now
interpretable as a shell boundary effect. The ring sits at the radius
where Quaoar's local fold field and the solar system's external field
are in balance. Inside this radius: material is in Quaoar's shell,
experiences Quaoar's $B$, accretes normally. Outside this radius:
material is in the solar system shell, subject to tidal disruption.

The ring persists at the boundary because:
\begin{enumerate}
  \item $g_{\mathrm{local}} \approx g_{\mathrm{external}}$ --- neither
        dominates, so neither accretes nor disrupts completely.
  \item $B_{\mathrm{Quaoar}} \approx B_{\mathrm{solar}}$ at the boundary
        --- the $B$-differential is near zero, reducing the effective
        tidal pressure.
\end{enumerate}
This is a testable prediction: ring systems should preferentially form
near the Hill sphere boundary, not at the classical Roche limit.
The classical Roche limit assumes a pure $1/r^2$ field without the
$B$ term; the fold Roche limit includes the $B$-differential and
predicts a shifted stability boundary.

\subsection{Connection to the Fold Scaling Kernel}

A companion paper~\cite{quaoar} derives the fold ring stability
condition from the lock condition $T_{\mathrm{fold}}(r) =
T_{\mathrm{orb}}(r)$. This yields the renormalization relation:
\begin{equation}
  \left(\frac{r}{l_P}\right)^n = \frac{\kappa_{\mathrm{eff}}}{\kappa_{\mathrm{sim}}},
  \label{eq:kernel_mod3}
\end{equation}
where $\kappa_{\mathrm{eff}} = 2\pi v_{\mathrm{orb}}/c$ and
$\kappa_{\mathrm{sim}} = 2.0$ is the Planck-scale simulation threshold.
This relation describes precisely how the fold threshold $\kappa$
is renormalized as you scale from the Planck length to macroscopic
ring radii.

This is the Module~3 shell hierarchy in quantitative form:
the effective fold threshold at any scale is set by the ratio
$(r/l_P)^n$, where $n$ is measured from the ring system's
body type (giant planet: $n \approx -0.093$; TNO: $n \approx -0.135$;
Centaur: $n \approx -0.160$).
Seven ring systems across four body types follow this three-population
structure, providing the first observational constraint on how fold
dynamics scale from Planck to planetary scales.

The scale-dependent $B$ field of Module~3 and the scale-dependent
$\kappa_{\mathrm{eff}}$ of the ring kernel are two aspects of the
same underlying physics: fold dynamics that weaken with scale in a
structured, measurable way.

%──────────────────────────────────────────────────────────────────────
\section{Open Problems}

\textbf{Open Problem 1 (Fold Cosmology).}
The magnitude of $B_{\mathrm{boundary}}$ and its connection to the
observed dark energy density requires a fold cosmological model ---
expansion dynamics, scale factor evolution, and the T3 conservation
law applied at cosmological scales. This is the most important open
problem in the present work.

\textbf{Open Problem 2 (Scale Invariance Proof).}
Section~2.3 provides a nondimensionalization argument showing that
the shell boundary condition is genuinely scale-free (not merely
lacking explicit scale).
What remains unproven is whether the full dynamical solution ---
including initial conditions and boundary conditions at each scale ---
preserves this invariance.
A simulation confirming that Module~2.3 shell formation dynamics
produce the same qualitative behavior (same $C = 1.014$ plateau,
same transport suppression) at widely separated physical scales
(e.g.\ $r \sim l_P$ and $r \sim 10^{35} l_P$) would close this problem.

\textbf{Open Problem 3 (Dark Matter Sector).}
The two-sector fold field (visible $\kappa_v$ + dark $\kappa_d$) needs
simulation and analytical development. Key questions: what determines
$\kappa_d/\kappa_v$? How does the dark sector cluster differently?
Does the ratio of dark to visible matter density fall out of fold
dynamics, or is it a free parameter?

\textbf{Open Problem 4 (Quaoar Simulation).}
The shell boundary prediction for ring systems needs a simulation:
run fold diffusion around a point mass within an external field,
measure where the B-differential drops to near zero, and check
whether ring stability occurs at that radius rather than the
classical Roche limit.

%──────────────────────────────────────────────────────────────────────
\section{Summary}

The shell-formation dynamics of Module~1 extend naturally to all scales.
The shell boundary is the surface where local and external fold fields
balance (Definition~2.1). The background term $B$ screens completely
inside each shell, becoming the local reference. The cosmological shell
--- closed during the 2D$\to$3D transition --- has no external reference,
so its $B_{\mathrm{boundary}}$ is anchored at shell-closure conditions.

As structure forms, $B_{\mathrm{interior}}$ depletes. The growing
differential $\Delta B = B_{\mathrm{boundary}} - B_{\mathrm{interior}}$
acts as an outward pressure without requiring $B < 0$ --- only that
the interior is lower than the boundary, as in a closed pressurized
system. This is the fold dark energy candidate.

The dark sector arrives late because it cannot exist before shell closure.
The coincidence problem is resolved because dark energy is sourced by
the same structure formation that produces visible matter clustering.

\begin{center}
\fbox{
\begin{minipage}{0.85\textwidth}
\vspace{4pt}
\textbf{The fold shell hierarchy:}\\[4pt]
Shell boundary: $g_{\mathrm{local}}(r^*) = g_{\mathrm{ext}}(r^*)$
\hfill (bias to the larger)\\[2pt]
$B$-screening: $B_\ell = (M_{\mathrm{total},\ell} -
M_{\mathrm{structure},\ell})/V_\ell$
\hfill (T3 within shell)\\[2pt]
$\Delta B = B_{\mathrm{boundary}} - B_{\mathrm{interior}} \geq 0$
\hfill (grows as structure forms)\\[2pt]
$\Delta B \to B_{\mathrm{boundary}}$ as $a \to \infty$
\hfill ($w \to -1$ at late times)\\[4pt]
\textit{Dark energy is the $B$-differential of a closed fold system.}\\
\textit{It cannot appear before the shell closes.}
\vspace{4pt}
\end{minipage}
}
\end{center}

%──────────────────────────────────────────────────────────────────────
\section*{Acknowledgements}

The author designed and implemented all theoretical development and is
solely responsible for the scientific content.
Mathematical derivations and manuscript preparation were conducted with
AI assistance: Claude Sonnet~4.6 --- ``Green'' --- (Anthropic, 2026).
All errors remain the responsibility of the author.

Code and data available at:
\texttt{https://github.com/PrimeFoldTheory/foundations}

\begin{thebibliography}{9}

\bibitem{mod1}
S.~Sowden, \textit{Prime Fold Theory: Module~I --- The Prime Substrate
and Foundations of the Fold Ontology}, Prime Fold Theory Project,
2025--2026.

\bibitem{mod21}
S.~Sowden, \textit{Module~2.1: Cross-System Simulation of Fold-Time
Invariance}, Prime Fold Theory Project, August 2025.

\bibitem{mod22}
S.~Sowden, \textit{Module~2.2: Mathematical Well-Posedness of Fold
Dynamics}, Prime Fold Theory Project, September 2025.

\bibitem{mod23}
S.~Sowden, \textit{Module~2.3: Shell Formation from Fold Dynamics},
Prime Fold Theory Project, April 2026.

\bibitem{gauge}
S.~Sowden, \textit{Gauge Structure from Fold Cycle Dynamics},
Prime Fold Theory Project, March 2026.

\bibitem{attractor}
S.~Sowden, \textit{Attractor Structure of Minimal Fold Systems},
Prime Fold Theory Project, March 2026.

\bibitem{gravity}
S.~Sowden, \textit{Gravitational Inverse-Square Law from Fold
Conservation}, Prime Fold Theory Project, April 2026.

\bibitem{planck_tick}
S.~Sowden, \textit{The Planck Tick as a Derived Quantity in Prime
Fold Theory}, Prime Fold Theory Project, April 2026.

\bibitem{quaoar}
S.~Sowden, \textit{Quaoar's Anomalous Ring System as a Fold
Field Constraint}, Prime Fold Theory Project, April 2026.

\end{thebibliography}

\end{document}
